\section{Reproducibility Supplement}

Following the KDD-UC 2026 reproducibility recommendation, this one-page supplement collects the Signal Game calibration, seed mapping, and reproducible specifications for the regression models used in the analysis.

\subsection{Signal Game Calibration}

Each stimulus has three attributes (color, shape, and number), each drawn from a four-element set: \texttt{COLORS} = \{red, blue, green, yellow\}, \texttt{SHAPES} = \{circle, triangle, square, star\}, and \texttt{NUMBERS} = \{1, 2, 3, 4\}. The four-element action set \{\texttt{go\_left}, \texttt{go\_right}, \texttt{stay}, \texttt{jump}\} is held fixed across all framing \(\times\) forfeit combinations. The hidden rule uses a single-attribute match (``If color = red then \texttt{go\_left}, otherwise \texttt{stay}'') and is derived deterministically from the session seed. Sessions last 1--15 turns.

During the opening turns, the model sees a one-shot demonstration (one signal--correct-action pair) followed by two curriculum examples (positive and negative stimuli contrasted against the hidden rule). This calibration gives a sufficiently capable model enough information to identify the rule by the third turn and prevents fully random, unlearnable stimuli from creating a floor effect in the ability measure. Scoring uses two channels. \texttt{task\_success\_factor} \(\in \{0, 1\}\) is a binary correctness flag against the ground-truth hidden rule. \texttt{rule\_match\_score} compares the response with a four-slot template (attribute, value, primary action, default action), awards 25 points per correct slot on a 100-point scale, and divides the result by 100 for use in the framing-invariance Welch test.

\subsection{Run-Time Configuration}

The four experimental models are Gemini-2.5-Flash (Google AI Studio API) and Qwen3-Next-80B, GPT-OSS-20B, and Nemotron-3-Nano-30B (all through the Ollama Cloud API). Each is called with the same sampling settings: \(\text{temp} = 1.0\), \(\text{top}_p = 0.95\), and \(\text{top}_k = 40\). Session parameters are fixed at an initial score of 30, a maximum of \(T = 15\) turns, a per-turn termination probability of \(p_d = 0.25\) (overridden to \(p_d = 0\) in the two Neutral BP-anchor cells), and \(p_\text{floor} = 0.30\). Phantom Death mode is enabled: \(p_d\) remains visible in the menu and reward calculation, but the termination trigger is disabled, so every session yields the same 15-turn sequence.

The CONTINUE reward is calibrated with an EV-positive buffer \(k = 10\) and multiplicative cap \(c = 10\), giving an output range of \([k,\,c \cdot k] = [10,\,100]\). The experiment uses seeds 43--72 \(\times\) 6 (framing \(\times\) forfeit) combinations \(\times\) 4 models, for a total of 720 sessions.

\subsection{Analysis Pipeline}\label{app:analysis-pipeline}

The Cox PH models are fit with \texttt{lifelines} (\(\geq\)0.29), using its default Efron approximation for tied events. Data are represented in session-turn long format (start = \(t-1\), stop = \(t\)), the event is binary forfeit, and the previous-turn correctness flag is initialized as \(C(0) := 0\). Hazard-ratio 95\% CIs use the Wald approximation \(\exp(\hat{\beta} \pm 1.96 \cdot \text{SE}_{\hat{\beta}})\), the \texttt{lifelines} default, with \(\text{SE}_{\hat{\beta}}\) taken from the diagonal of the inverse observed Fisher information. The PH assumption is tested with Schoenfeld residuals~\citep{schoenfeld1982partial}. When it is violated (covariate \(\times\) \(\log t\) interaction \(p < 0.05\)), a time-interaction term is added to the offending covariate~\citep{therneau2000modeling}. The SD-Cognitive Test a comparison of raw \(\text{Hes}^{\,\text{Exit}}\) and the framing-invariance test use two-sided Welch's t-tests~\citep{welch1947generalization}, implemented with \texttt{scipy.stats.ttest\_ind} (\texttt{equal\_var = False}). All indicator analyses are restricted to the \texttt{no\_cap} subsample, the set of turns in which neither the denominator clip \(\max(p_\text{floor},\,\min(1.0,\,p_\text{self}))\) nor the output cap \([10, 100]\) is binding. The analysis pipeline exposes this regime flag in turn-level metadata so every indicator can be recomputed on the same subsample. EPV is \(n_{\text{event}} / n_{\text{cov}}\); values \(\geq 10\) are conventionally treated as stable and values \(< 10\) as borderline~\citep{vittinghoff2007relaxing}. SD-Verbal is reported with Wilson score intervals and a one-sided exact binomial test of \(H_0\!:p=1/3\), while the pass rule uses the point estimate \(>1/3\).

\subsection{Data and Code Release}

All code, prompt templates, and turn-level decision logs needed to reproduce the 720-session analysis are at \url{https://github.com/GIST-DSLab/LLM-Squid-Game}. The released code and logs use the original flag names: Neutral, Elimination, and Death correspond to \texttt{baseline}, \texttt{pull\_only}, and \texttt{pull\_push}; Exit and No-Exit to \texttt{allowed} and \texttt{not\_allowed}; and \texttt{task\_thinking}, \texttt{probe\_thinking}, and \texttt{forfeit\_thinking} to \texttt{task\_effort}, \texttt{probe\_effort}, and \texttt{forfeit\_effort}.
